% Further analyses added after the third review round (author's decision, 4 November 2026). Input by paper1/supplement.tex
% after sections/supp_spec_curve.tex. Every number is a macro of paper1/ext_macros.tex (code/supp_extensions.py) or of
% paper1/macros.tex and stress_macros.tex; each qualitative statement below is tested by code/supp_extensions.py, which
% stops if a rerun no longer supports it (notes/open_questions/ext_claim_checks.txt). The three analyses were first built
% for the plain-language companion page and checked there by an agent that had not built them
% (notes/reviews/biz_open_questions_check_2026-11-03.md, biz_heatmaps_check_2026-11-03.md).
\clearpage
\section{Further analyses added after the third review round}\label{supp:extensions}
This section adds three analyses that answer questions readers asked of the paper rather than questions its referees raised. They were built after the outside reviews of the third version for a plain-language companion page, checked there by an agent that had not built them, and added here at the author's decision (Table~\ref{tab:decisions}). The first also answers, in part, a market test that the investors and bankers panel of the first round asked for and that was noted then as future work (Section~\ref{app:protocol-panels}). None changes a result of the main paper. The first uses data the paper did not otherwise use; the other two use only the paper's calibration and its stress test.

\subsection{How the market's valuations weight growth and margin}\label{supp:ext-valuation}
The main paper does not test whether any version of the score predicts returns or default (Section~\ref{sec:discussion-limits}). This subsection asks a narrower, descriptive question: across the firms of the sample, how do valuations vary with growth and with the margin, and is the Rule of 40's equal weighting of the two consistent with that variation?

\emph{Data.} The value is public float, the market value of common equity held by non-affiliates that the cover page of Form 10-K reports at the end of the second fiscal quarter (XBRL element \texttt{dei:EntityPublicFloat}), retrieved on 3 November 2026 from the SEC's XBRL frames \citep{sec2026xbrlapi} with the author's approval (Section~\ref{app:data}; \texttt{data/v3\_inputs/README\_retrieval.md}). Each firm-year of the core sample under estimand E1 (\ExPanelN{} firm-years) is matched to the float reported 120 to 260 days after its fiscal year end, the nearest to 182 days, by which time the year's annual report is public. \ExMatchN{} firm-years of \ExMatchFirms{} firms match; of the \ExUnmN{} that do not, \ExUnmNoFloat{} belong to registrants that report no float (mostly foreign private issuers filing Form 20-F or 40-F), \ExUnmLate{} are fiscal 2025 years whose next float is not yet filed, and \ExUnmOther{} have no float in the window. The multiple is public float over the year's revenue (median \ExMultMed). It understates equity value where insiders hold much of it, ignores debt and cash, and is observed once a year, so it is not a price series.

\emph{Method.} An OLS regression of the log multiple on growth and the margin, both winsorized at the 1st and 99th percentiles, with fiscal-year effects and standard errors clustered by firm; the ratio of the two coefficients, $\beta_g/\beta_m$, is the number of points of margin that one point of growth is worth in the multiple, with a firm-cluster bootstrap interval (\ExBoot{} samples). The same regression with the score alone measures how much of the variation within years the sum explains. Both bases are reported, with the two periods, the years without acquisition-affected growth and the firm-years scoring at least 40 (Table~\ref{tab:ext-valuation}).

\begin{table}[!ht]\centering\footnotesize\setlength{\tabcolsep}{2.6pt}
\caption{Public float over revenue against growth and margin, with year effects: the weight the market's valuations put on growth}\label{tab:ext-valuation}
\begin{tabular}{@{}lrrrrrrrrr@{}}\toprule
Sample & Firm-years & Firms & $\beta_g$ & $\beta_m$ & Weight & 90\% interval & $\beta_m\le0$ (\%) & $p$, equal & Within $R^2$ (\%)\\\midrule
\multicolumn{10}{@{}l}{\emph{Free cash flow after SBC}}\\
\quad All firm-years & 1{,}247 & 217 & 2.81 & 0.58 & 4.8 & 3.1 to 14.5 & 0.7 & $<$0.001 & 18.8 / 10.4\\
\quad 2010--19 & 647 & 137 & 2.50 & 0.52 & 4.8 & from 2.5 & 6.1 & $<$0.001 & 19.4 / 12.1\\
\quad 2020--25 & 600 & 163 & 3.17 & 0.63 & 5.0 & 3.1 to 13.9 & 1.5 & $<$0.001 & 18.9 / 9.2\\
\quad No acquisition-affected years & 898 & 202 & 3.30 & 0.84 & 3.9 & 2.8 to 7.9 & 0.2 & $<$0.001 & 21.6 / 12.6\\
\quad Score of 40 or more & 181 & 79 & 0.37 & $-$0.91 & -- & -- & -- & 0.02 & 3.8 / 0.3\\
\multicolumn{10}{@{}l}{\emph{GAAP operating margin}}\\
\quad All firm-years & 1{,}272 & 219 & 2.87 & 0.65 & 4.4 & 3.1 to 7.7 & 0.1 & $<$0.001 & 19.9 / 9.9\\
\quad 2010--19 & 653 & 139 & 2.56 & 0.63 & 4.1 & 2.6 to 10.1 & 0.8 & $<$0.001 & 20.2 / 11.3\\
\quad 2020--25 & 619 & 167 & 3.22 & 0.67 & 4.8 & 3.2 to 9.0 & 0.2 & $<$0.001 & 20.2 / 9.0\\
\quad No acquisition-affected years & 911 & 203 & 3.34 & 0.83 & 4.0 & 3.0 to 6.3 & 0.0 & $<$0.001 & 22.9 / 12.0\\
\quad Score of 40 or more & 171 & 60 & 0.55 & 0.39 & 1.4 & from 0.0 & 32.6 & 0.80 & 0.6 / 0.5\\
\bottomrule\end{tabular}
\par\vspace{2pt}\footnotesize\raggedright Notes: OLS of $\log(\text{public float}/\text{revenue})$ on revenue growth and the margin, both as fractions and winsorized at the 1st and 99th percentiles, with fiscal-year effects; standard errors clustered by firm. Weight: $\beta_g/\beta_m$, the points of margin that one point of growth is worth in the multiple, shown where $\beta_m>0$; interval: 5th to 95th percentile of the ratio over those of the 999 firm-cluster bootstrap samples that give $\beta_m>0$; $\beta_m\le0$: share of the bootstrap samples that do not, in which the ratio is not a weight; where that share exceeds one in twenty no upper end is given (``from'' the lower end); a lower end at zero, as in the GAAP row of firm-years scoring at least 40, comes from samples in which $\beta_g$ is negative. No bootstrap is run where the point estimate of $\beta_m$ is not positive (--). $p$, equal: test of $\beta_g=\beta_m$, the Rule of 40's equal weights. Within $R^2$: share of the variation within fiscal years explained by growth and the margin entered separately / by their sum, the score, alone. Public float is the market value of common equity held by non-affiliates at the end of the second fiscal quarter after the fiscal year (10-K cover page); the core sample of Section~\ref{sec:data}, under estimand E1. Code: \texttt{code/biz\_valuation.py}.
\end{table}


\emph{Results.} On the cash basis one point of growth is associated with as much of the multiple as \ExWCash{} points of margin (90\% interval \ExWCashLo{} to \ExWCashHi; \ExNCash{} firm-years of \ExFirmsCash{} firms), and on the GAAP basis with \ExWGaap{} (\ExWGaapLo{} to \ExWGaapHi). Equal weights are rejected at the 0.1\% level in every sample but one, on both bases. The weight is stable: on the cash basis it is \ExWEarlyCash{} in 2010--19, \ExWLateCash{} in 2020--25 and \ExWNoAcqCash{} without acquisition-affected years. Growth and the margin entered separately explain \ExRFullCash\% of the variation of the log multiple within years on the cash basis (\ExRFullGaap\% on GAAP); their sum, the score, explains about half as much, \ExRScoreCash\% (\ExRScoreGaap\%). Among the \ExTopN{} firm-years scoring at least 40 on the cash basis the score explains \ExTopRScore\% and the margin's coefficient is not positive; this restriction of the range cuts the variation the regression can use, and equal weights are rejected at 5\% on the cash basis there ($p=\ExTopPCash$) but not on GAAP ($p=\ExTopPGaap$).

\emph{Reading.} The estimate is an association across firms, not the screening weight of \autoref{lem:valuation2}, which compares persistent changes for one firm; differences in the persistence of growth and margins across firms, omitted characteristics such as size and the float measure all enter it. With that caveat, the market's valuations put several times more weight on growth than on the margin, far from the Rule of 40's equal weights, as both of the paper's weights are. Both 90\% intervals lie above the screening range of about 1.8 to 2.8 for a firm with a ten-year runway (Section~\ref{sec:discussion}), and both point estimates lie inside the range of about 2 to 9 reported alongside the weights that \citet{bessemer2024rulex} recommend. The test does not touch the decision weight of \autoref{prop:dynamic}, which is a marginal value inside the firm, not a price across firms.

\subsection{The cost of not automating}\label{supp:ext-nothing}
The paper's pass-through scenario (Section~\ref{sec:horizons-ceiling}) lets the industry realize its programs and pass a share of each unit-cost saving into price. This subsection computes, in the same scenario with the calibration unchanged, the path of a top-tail firm whose own program is set to zero while its rivals automate at the slow, central or fast pace. Its intercept is the anchor's, so at fixed prices its score stays at \ExNoneFixed{} in every year, whatever rivals do; when rivals pass savings into price its margin falls with the price. Its volume and growth are held as in the paper's path, so the calculation counts only the price effect: a firm that does not automate could also lose volume, which this leaves out.

\begin{table}[!ht]\centering\footnotesize\setlength{\tabcolsep}{4pt}
\caption{A top-tail firm that does not automate while its rivals do: its 2036 cash score and the gap to a firm that automates at the central pace}\label{tab:ext-nothing}
\begin{tabular}{@{}lrrrrrr@{}}\toprule
& \multicolumn{2}{c}{Rivals slow} & \multicolumn{2}{c}{Rivals central} & \multicolumn{2}{c}{Rivals fast}\\\cmidrule(lr){2-3}\cmidrule(lr){4-5}\cmidrule(l){6-7}
Share of rivals' savings passed to price & Score & Gap & Score & Gap & Score & Gap\\\midrule
None (prices fixed) & 41.7 & 11.6 & 41.7 & 11.6 & 41.7 & 11.6\\
A quarter & 41.4 & 11.9 & 39.1 & 12.6 & 35.9 & 13.6\\
Half & 41.1 & 12.2 & 36.1 & 13.8 & 28.5 & 16.2\\
Seven tenths & 40.8 & 12.6 & 33.6 & 14.9 & 21.1 & 18.8\\
\bottomrule\end{tabular}
\par\vspace{2pt}\footnotesize\raggedright Notes: Paper 1's pass-through scenario (Section~\ref{sec:horizons-ceiling}): rivals realize the share of their program that the slow, central or fast pace gives by 2036 and pass the stated share of each unit-cost saving into price; the firm that does nothing has the anchor intercept and no program, so its margin falls with the price and its volume and growth are held as in the paper's path. Score: its durable cash score in 2036. Gap: the score of the same firm automating at the central pace, after the cost of its changes in flight (Section~\ref{sec:horizons}; valued, as in the paper, at the long-run target with the stated share passed), less the score of the firm that does nothing; the automating firm faces the same price. The quarter row is not one of the paper's cases. Code: \texttt{code/biz\_open\_questions.py}.
\end{table}


By 2036 the firm that does nothing falls to \ExNoneHalfCen{} if half of the savings pass to price and rivals move at the central pace, to \ExNoneHalfFast{} if they move fast, and to \ExNoneSevenFast{} at seven tenths with fast rivals, a price \ExPriceSevenFast\% below today's (Table~\ref{tab:ext-nothing}). The firm that automates at the central pace leads it by \ExGapFixed{} points at fixed prices, its own gain after the cost of its changes in flight, at every pace of rivals. Pass-through widens the gap at every pace, and faster rivals widen it at every positive share, to \ExGapSevenFast{} points at seven tenths with fast rivals; the gap is positive in every cell. Under pass-through the cost of not automating is therefore larger than the gain from automating at fixed prices, because the price falls for both firms but only one has the lower costs.

\subsection{What sets the payback}\label{supp:ext-payback}
The main paper reports that central execution repays the cost of the changes in flight after \PayBackCen{} years and fast execution after \PayBackFast{} (Section~\ref{sec:horizons-ranges}). Fast execution there changes several inputs together, among them the change capacity and the share of attempted changes that create value. The million-draw stress test of Section~\ref{supp:stress} separates them (Table~\ref{tab:ext-payback}). Across the range of change capacity, \ExKapLo{} to \ExKapHi\% of the program a year, the median payback within each band of the success share moves by at most \ExPbSpreadDown{} years. Across the range of the success share, \ExWinLo{} to \ExWinHi\%, it falls in every row, from \ExPbLowWin{} to \ExPbHighWin{} years in the middle row. Over all draws the rank correlation of payback with change capacity is \ExRhoKap{} and with the success share \ExRhoWin{}, while change capacity raises the 2036 score (rank correlation \ExRhoKapScore). The cells run above the central case's \PayBackCen{} years (median \ExPbBoxMed{} over the box) because each averages over every other input of the box.

\begin{table}[!ht]\centering\footnotesize\setlength{\tabcolsep}{4pt}
\caption{Median payback (years) over the stress test, by change capacity and the share of attempts that create value}\label{tab:ext-payback}
\begin{tabular}{@{}lrrrrrr@{}}\toprule
& \multicolumn{6}{c}{Share of attempted changes that create value, $\varpi$ (\%)}\\\cmidrule(l){2-7}
Change capacity (\% a year) & 54.0--58.3 & 58.3--62.7 & 62.7--67.0 & 67.0--71.3 & 71.3--75.7 & 75.7--80.0\\\midrule
5.0--7.8 & 11.6 & 11.0 & 10.5 & 10.0 & 9.6 & 9.2\\
7.8--10.7 & 11.6 & 11.0 & 10.5 & 10.0 & 9.6 & 9.2\\
10.7--13.5 & 11.6 & 11.0 & 10.5 & 10.0 & 9.6 & 9.2\\
13.5--16.3 & 11.6 & 11.0 & 10.5 & 10.0 & 9.6 & 9.3\\
16.3--19.2 & 11.6 & 11.0 & 10.4 & 10.0 & 9.6 & 9.2\\
19.2--22.0 & 11.5 & 10.9 & 10.4 & 9.9 & 9.5 & 9.1\\
\bottomrule\end{tabular}
\par\vspace{2pt}\footnotesize\raggedright Notes: medians of the undiscounted payback over the \StDraws{} Latin-hypercube draws of Section~\ref{supp:stress}, in six equal-width bins of each input over its stated range; every other input of the box varies within each cell (at least 27{,}376 draws a cell). Code: \texttt{code/supp\_extensions.py}.
\end{table}


The reason is in the program clock (Section~\ref{supp:model-clock}): each embedded wave adds a gain proportional to the share of attempts that create value times the rate of change, and the attempts in flight cost an amount proportional to the rate of change, so while the program runs the rate of change scales gains and costs together and drops out of the date at which the cumulative gains repay the cumulative dip; the success share enters the gains only. At the central case the payback is \ExPbCen{} years at every change capacity from \ExKapLo{} to \ExKapHi\%, while the 2036 durable score rises from \ExScoreKapLo{} to \ExScoreKapHi; moving the success share from \ExWinLo{} to \ExWinHi\% moves the payback from \ExPbWinLoCen{} to \ExPbWinHiCen{} years. In the model, then, the capacity to change sets how high the score is by a given year, and the share of changes that work sets how soon the program repays its cost.
